\chapter{Metodologia}

Neste capítulo, são apresentados os métodos utilizados para o desenvolvimento da análise multivariada, organizados segundo as etapas do processo KDD (\textit{Knowledge Discovery in Databases}): Limpeza de dados, Integração de dados, Seleção de dados, Transformação de dados, Mineração de dados, Avaliação de padrões e Apresentação do conhecimento.

% ============================================================
\section{Delimitação do estudo}
% ============================================================

\subsection{Unidade de análise}
O Brasil é um país com 5.570 municípios, 26 estados e o Distrito Federal, com a população estimada em 
203.080.756, segundo o censo de 2022. O estado de Minas Gerais situa-se na região sudeste do Brasil e possui 853 municípios. Sua população em 2022 era de 
20.539.989 pessoas habitantes, destes, cerca de 10.524.280 milhões de mulheres. 

A unidade principal de análise adotada neste trabalho é o município-ano, ou seja, cada observação corresponde a um dos 853 municípios de Minas Gerais, identificado pelo código do Índice Brasileiro de Geografia e Estatística (IBGE), em um determinado ano. 

\subsection{População-alvo}

O estudo é restrito a mulheres residentes no estado de Minas Gerais. A faixa etária considerada abrange de 0 até 75 anos ou mais, seguindo a estrutura presente nas bases de dados.

\subsection{Períodos de análise e bases utilizadas}

Em razão da heterogeneidade entre as fontes de dados, cada base apresenta uma cobertura temporal distinta:

\begin{enumerate}
    \item \textbf{Taxa de incidência de câncer de mama} (RHC/INCA): 1985 a 2023;
    \item \textbf{Taxa de mortalidade por câncer de mama} (INCA/GOV): 1979 a 2024;
    \item \textbf{Índice de Desenvolvimento Humano Municipal -- IDHM} (IPEA): 2013 a 2024;
    \item \textbf{População por município} (IBGE): anos censitários de 1970, 1980, 1991, 2000, 2010 e 2022;
    \item \textbf{Hospitais habilitados em oncologia} (GOV): situação em 2025.
\end{enumerate}

A janela temporal de integração final, ou seja, o período em que todas as bases se sobrepõem com dados disponíveis, foi definida considerando a interseção entre as fontes. As diferenças de cobertura histórica refletem limitações institucionais de cada sistema de informação: enquanto o SIM registra óbitos desde 1979, o RHC passou a consolidar dados de incidência de forma sistemática a partir dos anos 1985. Já os dados censitários seguem outro padrão de periodicidade, o que impõe uma granularidade diferente das demais séries anuais. Essas diferenças são levadas em conta nas etapas de transformação e integração descritas adiante.

% ============================================================
\section{Limpeza de dados}
% ============================================================

\subsection{Critérios de inclusão}

Foram incluídos apenas registros que atendessem simultaneamente às seguintes condições:

\begin{enumerate}
    \item Município de residência localizado no estado de Minas Gerais, identificado pelo código IBGE;
    \item Sexo feminino;
    \item Para os dados relacionados ao câncer de mama, registros com classificação de neoplasia maligna da mama (CID-10 C50);
    % \item Presença das variáveis essenciais para a análise multivariada (município, ano e atributo de interesse).
\end{enumerate}

\subsection{Critérios de exclusão}

Foram excluídos registros que apresentassem:

\begin{enumerate}
    % \item Código de município ausente, inválido, não padronizável ou fora do conjunto oficial de municípios de Minas Gerais;
    % \item Registros fora do recorte do estudo (sexo diferente de feminino, quando aplicável; ou classificação diagnóstica diferente de CID-10 C50 para a base RHC/INCA);
    \item Registros com ano fora da janela temporal definida para cada fonte;
    \item Duplicatas exatas após padronização das chaves de identificação;
    \item Registros com inconsistências não recuperáveis, como campos estruturais vazios (município/ano) e idades ignoradas.
\end{enumerate}

\subsection{Ausência de dados}

Valores em branco ou ausentes foram mantidos como estão, sem imputação, e identificados como vazio nas tabelas geradas. 

\subsection{Duplicidades, inconsistências e padronização}

Foram aplicadas os seguintes tratamentos:

\begin{enumerate}
    \item Normalização textual de municípios, regionais e hospitais (remoção de acentos e padronização de caixa para minúsculas);
    \item Padronização dos códigos de município pelo padrão IBGE de sete dígitos;
    \item Remoção de duplicatas exatas identificadas por combinação de chave município-ano-variável e consolidação por soma quando necessário em tabelas agregada;
    \item Tratamento de inconsistências estruturais nas planilhas de entrada como a remoção de linhas totalmente vazias e ajuste de nomes de colunas;
    \item Harmonização de rótulos categóricos (faixa etária, escolaridade e raça/cor) e padronização de colunas temporais no formato \texttt{ANO\_CATEGORIA};
    \item Conversão padronizada de campos numéricos (incluindo formato brasileiro com ponto de milhar e vírgula decimal), com manutenção explícita de ausências quando não há valor observável;
    \item Organização estrutural das tabelas em formato matriz com municípios nas linhas (identificados pelos campos \texttt{municipio} e \texttt{cod\_ibge}) e variáveis nas colunas, precedidas por identificadores temporais (anos) e seguidas por atributos específicos;
    \item Padronização das faixas etárias em categorias de dez anos de amplitude (\texttt{0 a 9}, \texttt{10 a 19}, \ldots, \texttt{75 ou mais}), mantendo coerência com a estrutura presente nas fontes originais.
\end{enumerate}

% \subsection{Considerações sobre viés}

% Algumas limitações inerentes às fontes de dados podem introduzir viés na análise. Entre elas, destacam-se a subnotificação histórica de casos, especialmente em municípios menores e em períodos anteriores à expansão do SUS; as mudanças de classificação diagnóstica ao longo do tempo, como as revisões da CID; e a cobertura desigual dos sistemas de informação entre regiões do estado. Essas questões são discutidas com mais detalhe na seção de Limitações.

% ============================================================
\section{Integração de dados}
% ============================================================

\subsection{Fontes de dados}

As bases utilizadas foram obtidas junto ao INCA (Instituto Nacional de Câncer), ao Governo do Estado de Minas Gerais (GOV-MG), ao IPEA (Instituto de Pesquisa Econômica Aplicada) e ao IBGE (Instituto Brasileiro de Geografia e Estatística), todas de acesso público. O IFDM (Índice FIRJAN de Desenvolvimento Municipal) e o Integrador RHC (Registro Hospitalar do Câncer) também foram utilizados como fontes complementares.

\subsection{Chave de junção}

A integração entre as bases foi realizada utilizando o código IBGE do município como chave primária. Antes da junção, todos os campos de código e nome de município passaram por padronização conforme descrito na etapa de limpeza. 

\subsection{Estratégia de integração}

A junção foi realizada de forma incremental, adicionando progressivamente cada base ao conjunto central. A estratégia de junção adotada para a construção do painel final foi do tipo \textit{inner join} por município, de modo a garantir que apenas municípios com dados disponíveis integrassem o conjunto consolidado.

\subsection{Validação pós-integração}

Após a integração, foi verificada a cobertura final por variável, incluindo o número de municípios válidos, o número de anos cobertos e o percentual de preenchimento por coluna. 

% Eventuais perdas relevantes de cobertura foram registradas e consideradas na interpretação dos resultados.

% ============================================================
\section{Seleção de dados}
% ============================================================

A partir da análise exploratória das bases disponíveis, foram definidos os atributos de interesse com base em dois critérios: relevância para a análise multivariada do câncer de mama e disponibilidade de série temporal suficiente para comparações municipais. Os atributos selecionados foram organizados em dois grupos:

\subsection{Atributos demográficos}

\begin{enumerate}
    \item Quantidade de hospitais habilitados em oncologia em Minas Gerais, com situação referente a 2025;
    \item Taxa de incidência de câncer de mama por município, com abertura por ano de diagnóstico, escolaridade, faixa etária e raça/cor, no período de 1985 a 2023;
    \item População feminina por município, para a faixa etária de 0 até 75 anos ou mais, nos anos censitários de 1970, 1980, 1991, 2000, 2010 e 2022.
\end{enumerate}

\subsection{Atributos socioeconômicos}

\begin{enumerate}
    \item Taxa de mortalidade por câncer de mama por município, no período de 1979 a 2024;
    \item Índice de Desenvolvimento Humano Municipal (IDHM) -- nas dimensões geral, educação, saúde e emprego e renda -- por município, de 2013 a 2024.
\end{enumerate}

\subsection{Dicionário de variáveis}

Para garantir rastreabilidade e padronização semântica na etapa de integração, foi construído um dicionário de variáveis para cada tabela analítica. A descrição segue o padrão nome da tabela + lista de atributos, conforme o exemplo abaixo.

\textbf{incidencia\_municipio\_ano:} armazena indicadores de incidência de câncer de mama por município e ano em Minas Gerais. Seus atributos são:

\begin{itemize}
    \item \texttt{cod\_ibge}: código IBGE do município (7 dígitos);
    \item \texttt{municipio}: nome do município;
    \item \texttt{ano\_diagnostico}: ano de referência do diagnóstico;
    \item \texttt{faixa\_etaria}: categoria etária da paciente;
    \item \texttt{escolaridade}: nível de escolaridade registrado;
    \item \texttt{raca\_cor}: categoria de raça/cor autodeclarada;
    \item \texttt{casos\_c50}: número absoluto de casos de neoplasia maligna da mama (CID-10 C50);
    \item \texttt{pop\_feminina}: população feminina do município no ano de referência;
    \item \texttt{taxa\_incidencia\_100k}: taxa de incidência por 100 mil mulheres.
\end{itemize}

Como complemento ao texto descritivo, a Tabela~\ref{tab:exemplo_dicionario_variaveis} apresenta um modelo de documentação tabular do dicionário de variáveis.

\begin{table}[htbp]
\centering
\caption{Exemplo de dicionário de variáveis da tabela \texttt{incidencia\_municipio\_ano}}
\label{tab:exemplo_dicionario_variaveis}
\small
\begin{tabular}{|l|p{4.8cm}|l|p{3.5cm}|c|}
\hline
\textbf{Variável} & \textbf{Descrição} & \textbf{Tipo} & \textbf{Formato/Valores} & \textbf{Nulo} \\
\hline
\texttt{cod\_ibge} & Código IBGE do município & Inteiro & 7 dígitos & Não \\
\hline
\texttt{municipio} & Nome do município & Texto & String & Não \\
\hline
\texttt{ano\_diagnostico} & Ano de referência do diagnóstico & Inteiro & 1985--2023 & Não \\
\hline
\texttt{faixa\_etaria} & Faixa etária da paciente & Categórico & Faixas padronizadas & Sim \\
\hline
\texttt{escolaridade} & Escolaridade registrada & Categórico & Categorias RHC/INCA & Sim \\
\hline
\texttt{raca\_cor} & Raça/cor autodeclarada & Categórico & Categorias padronizadas & Sim \\
\hline
\texttt{casos\_c50} & Número absoluto de casos (CID-10 C50) & Inteiro & $\geq 0$ & Não \\
\hline
\texttt{pop\_feminina} & População feminina municipal & Inteiro & $\geq 0$ & Sim \\
\hline
\texttt{taxa\_incidencia\_100k} & Taxa por 100 mil mulheres & Decimal & $\geq 0$ & Sim \\
\hline
\end{tabular}
\end{table}

% ============================================================
\section{Transformação de dados}
% ============================================================

\subsection{Coleta automatizada}

A coleta dos dados foi realizada com o uso da linguagem Python, versão 3.12, por meio de técnicas de \textit{web scraping} para extração e automatização do processo. Foram utilizadas as bibliotecas Selenium para acessar as páginas de interesse, extrair os dados relevantes e armazená-los em formato CSV para análise posterior.

\subsection{Variáveis derivadas}

A partir dos dados brutos coletados, foram construídas as seguintes variáveis:

\begin{enumerate}
    % \item \textbf{Taxa de incidência e mortalidade por 100 mil habitantes:} calculada a partir do número absoluto de casos e óbitos dividido pela população municipal estimada, multiplicado por 100.000. Essa transformação permite comparar municípios de diferentes tamanhos populacionais;
    \item \textbf{Categorias do IDHM:} os valores contínuos do índice foram classificados em quatro faixas, conforme a seguinte escala:
    \begin{enumerate}
        \item Desenvolvimento Crítico: IDHM inferior a 0,400;
        \item Desenvolvimento Baixo: IDHM entre 0,400 e 0,599;
        \item Desenvolvimento Moderado: IDHM entre 0,600 e 0,799;
        \item Desenvolvimento Alto: IDHM igual ou superior a 0,800.
    \end{enumerate}
\end{enumerate}

\subsection{Padronização temporal}


\subsection{Tabelas consolidadas}

Como resultado da etapa de transformação, foram construídas três tabelas principais:

\begin{enumerate}
    \item \textbf{Tabela 1 -- Mortalidade por câncer de mama:} número de óbitos por município em mulheres, de 1979 a 2024, em que cada linha representa um município e as colunas correspondem aos anos, com o valor agregado bruto de óbitos;
    \item \textbf{Tabela 2 -- IDHM por município:} índice de desenvolvimento humano municipal (geral, educação, saúde e emprego e renda), de 2013 a 2024, em que as linhas representam os municípios e as colunas os anos. A ausência de dados é identificada com asterisco (*);
    \item \textbf{Tabela 3 -- Incidência de câncer de mama:} taxa de incidência por município, com abertura por ano de diagnóstico, escolaridade, faixa etária e raça/cor, no período de 1985 a 2023. A essa tabela foram agregadas as informações sobre população municipal e hospitais habilitados.
    \item \textbf{Tabela 4 -- População por município:} população feminina por município, para a faixa etária de 0 até 75 anos ou mais, nos anos censitários de 1970, 1980, 1991, 2000, 2010 e 2022.
\end{enumerate}

% ============================================================
\section{Mineração de dados}
% ============================================================

% TODO: detalhar algoritmos e métricas de avaliação

% A etapa de mineração de dados tem como objetivo identificar padrões e estruturas nos dados consolidados que não seriam perceptíveis por inspeção direta. As perguntas analíticas que orientam essa etapa são:

% \begin{enumerate}
%     \item Quais municípios apresentam padrões semelhantes de incidência, mortalidade e condições socioeconômicas ao longo do tempo?
%     \item Existem agrupamentos regionais com características epidemiológicas distintas?
%     \item Como as variáveis socioeconômicas se associam à evolução temporal da incidência e da mortalidade?
% \end{enumerate}

% Para responder a essas perguntas, serão aplicadas técnicas de agrupamento (\textit{clustering}), análise de tendência temporal e visualização de padrões regionais. Os algoritmos específicos e as métricas de avaliação serão detalhados após a definição final do pipeline de análise.

% ============================================================
\section{Avaliação de padrões}
% ============================================================

% Os padrões identificados na etapa de mineração serão avaliados segundo três critérios principais:

% \begin{enumerate}
%     \item \textbf{Relevância substantiva:} os padrões devem ser interpretáveis à luz do conhecimento epidemiológico e úteis para subsidiar decisões de saúde pública;
%     \item \textbf{Coerência epidemiológica:} os agrupamentos e associações encontrados devem ser consistentes com evidências da literatura sobre fatores de risco e acesso ao diagnóstico no câncer de mama;
%     \item \textbf{Robustez temporal:} os padrões devem apresentar estabilidade ao longo dos anos analisados, e variações abruptas devem ser investigadas quanto a possíveis artefatos das fontes de dados.
% \end{enumerate}

% O objetivo final dessa etapa é encontrar o perfil epidemiológico do câncer de mama em Minas Gerais, identificando as regiões mais vulneráveis, os fatores de risco associados e as possíveis lacunas no acesso ao diagnóstico e tratamento. A partir disso, será possível subsidiar políticas públicas e estratégias de prevenção direcionadas às populações mais vulneráveis, visando a redução da incidência e mortalidade pela doença.

% ============================================================
\section{Apresentação do conhecimento}
% ============================================================

% Os resultados serão apresentados por meio de diferentes formatos visuais, cada um voltado a responder perguntas específicas da análise:

% \begin{enumerate}
%     \item \textbf{Mapas de calor municipais:} para visualizar a distribuição espacial da incidência, mortalidade e IDHM ao longo do território mineiro, evidenciando desigualdades regionais;
%     \item \textbf{Séries temporais:} para acompanhar a evolução dos indicadores ao longo dos anos, permitindo identificar tendências e rupturas;
%     \item \textbf{Tabelas-síntese:} para consolidar os principais achados por grupo de municípios ou região, facilitando a comparação entre perfis distintos.
% \end{enumerate}

% Em conjunto, esses elementos visam não apenas descrever o cenário atual, mas oferecer subsídios analíticos concretos para o planejamento em saúde e a formulação de políticas públicas mais direcionadas às realidades locais de Minas Gerais.

% ============================================================
\section{Reprodutibilidade}
% ============================================================

Todo o pipeline de análise foi desenvolvido em Python 3.12, com registro das versões das bibliotecas utilizadas e da data de extração de cada base de dados. Os scripts estão organizados em ordem de execução e documentados de forma a permitir a reprodução integral do processo. As fontes de dados são todas públicas e de acesso aberto, com os endereços de coleta registrados na documentação do projeto.

% ============================================================
\section{Limitações}
% ============================================================

Esta análise apresenta algumas limitações que devem ser consideradas na interpretação dos resultados:

\begin{enumerate}
    \item \textbf{Subnotificação:} é esperada variação na completude dos registros, especialmente em municípios menores e em períodos históricos anteriores à expansão dos sistemas de informação em saúde;
    \item \textbf{Mudanças metodológicas entre fontes e anos:} revisões na Classificação Internacional de Doenças (CID) e alterações nos critérios de registro ao longo do tempo podem introduzir inconsistências na série histórica;
    \item \textbf{Ausência de variáveis clínicas:} aspectos como estágio ao diagnóstico, tipo histológico e tratamento recebido não estão disponíveis de forma sistemática nas bases públicas utilizadas, o que limita a profundidade da análise clínica;
    \item \textbf{Dados censitários:} a estimativa populacional por município baseia-se em anos censitários, o que introduz imprecisão nos denominadores das taxas para anos intercensitários.
\end{enumerate}

Essas limitações não invalidam os achados, mas orientam a interpretação com cautela, especialmente nas comparações entre municípios com históricos de cobertura muito distintos.

% ============================================================
\section{Aspectos éticos e LGPD}
% ============================================================

Este trabalho utiliza exclusivamente dados públicos e agregados, disponibilizados por instituições governamentais e de pesquisa. Não há acesso a informações individuais identificáveis, e todos os dados são tratados em nível municipal. A pesquisa tem finalidade estritamente acadêmica e de interesse público, voltada à compreensão de desigualdades em saúde e ao subsídio de políticas públicas, estando em conformidade com os princípios da Lei Geral de Proteção de Dados (LGPD).

